\section*{Discussion}

\subsection*{Model selection and the genetic architecture of each drug}

The consistent ranking of Random Forest above gradient-boosted trees, linear SVM, and regularized logistic regression across all four drugs and both the per-drug and multi-label formulations suggests that, on this feature representation, the bagging-based ensemble is more robust than boosting or linear alternatives to the high dimensionality (2{,}693 features), the correlation structure among nearby SNP positions, and the class imbalance of the dataset. The margins are narrow, however -- LightGBM trails Random Forest by 0.0015--0.0132 CV AUC-ROC depending on the drug (Table~\ref{tab:modelcomparison}) -- and with a single train/test split per benchmarking run we cannot claim these differences are statistically meaningful. The honest reading is that the tree ensembles are indistinguishable on this data and the linear models are worse.

The performance ordering across drugs is more informative than the ordering across models. Rifampicin and isoniazid resistance are each driven predominantly by a small number of high-penetrance mutations in \emph{rpoB} and \emph{katG} respectively \cite{who2023catalogue,walker2015whole}, which a classifier can exploit almost as a near-deterministic rule. Pyrazinamide is the opposite case: resistance is conferred by loss of \emph{pncA} function, and the substitutions that cause it are distributed across the entire length of the gene without hotspots -- more than 300 resistance-conferring substitutions have been characterized, many acting by reducing PncA protein abundance rather than by altering a catalytic residue \cite{yadon2017pnca}. A feature set of individual nucleotide positions is structurally ill-suited to that architecture: each causal variant is individually rare, so no single column carries much signal. This is a limitation of the representation, not of the classifier, and it bounds what any position-encoded model can achieve for this drug.

\subsection*{Performance in context: this pipeline is not state of the art}

The per-drug AUC-ROC values reported here (0.883--0.976) are at or below what established tools already achieve. TB-Profiler \cite{phelan2019tbprofiler} and the machine-learning-based GenTB \cite{groschel2021gentb} are both mature, widely deployed predictors covering 13 drugs; GenTB's own benchmark against 20{,}408 phenotyped isolates reports mean sensitivities of 77.6\% (GenTB random forest), 75.4\% (GenTB wide-and-deep network), 74.4\% (TB-Profiler) and 71.9\% (Mykrobe) across nine shared drugs, at specificities above 96\% \cite{groschel2021gentb}. For the multi-label problem specifically, DeepAMR \cite{yang2019deepamr} -- a multi-task denoising autoencoder built for co-occurrent resistance, the same problem our \texttt{MultiOutputClassifier} extension addresses -- reports AUCs of 0.944--0.987 on the first-line drugs, comfortably above our macro AUC-ROC of 0.919 and against a far more demanding target.

We therefore do not claim competitive accuracy, and this manuscript should not be read as proposing a replacement for those tools. Two caveats cut in both directions. First, cross-study comparison of these numbers is not sound: cohorts, label provenance, critical concentrations, train/test partitioning and drug coverage all differ, and benchmarking work has shown that reported performance for TB resistance prediction is strongly sensitive to which database and which isolate set is used \cite{ngo2019benchmarking}. Our own numbers are no more comparable to DeepAMR's than DeepAMR's are to TB-Profiler's. Second, that non-comparability is itself an indictment of the field rather than a defence of our results: without a shared, versioned benchmark it is not possible to say whether any of these methods is actually better than another. The contribution we can defend is the openly released feature matrix, code, models and per-isolate explanation layer -- infrastructure that makes such a comparison easier -- not the accuracy of the classifiers built on top of it.

\subsection*{Label quality bounds what can be learned}

A model can be no better than its labels, and ours are heterogeneous: phenotypes were aggregated across many source studies with unrecorded testing methods and critical concentrations. This matters most for pyrazinamide. Phenotypic pyrazinamide susceptibility testing in the Bactec MGIT 960 system is well documented to produce false-resistant results, driven by inoculum density and medium acidification, with reported false-resistance rates in the tens of percent and substantial improvement when inoculum is reduced \cite{mok2021pza,mustazzolu2019revisiting}. Some fraction of the pyrazinamide labels used here is therefore likely to be wrong, and in a specific direction: susceptible isolates mislabelled as resistant. Our weaker pyrazinamide performance consequently confounds three distinct effects -- a genuinely harder genetic architecture, a representation unsuited to it, and noisy ground truth -- and the present design cannot separate them. Any future evaluation of this pipeline on pyrazinamide should use a cohort with documented, quality-controlled DST, or quantitative MIC phenotypes of the kind assembled by the CRyPTIC compendium \cite{cryptic2022compendium}, rather than pooled binary calls.

\subsection*{Limitations of the SHAP analysis itself}

The interpretability layer, which is the paper's main claim to novelty, rests on assumptions this dataset violates. Shapley-value attributions are defined by averaging a feature's marginal contribution over subsets of the remaining features, and the standard estimators assume those features can be varied independently. Genomic features are the textbook case where that fails: linkage disequilibrium, the clonal population structure of the \emph{M. tuberculosis} complex, and the co-resistance correlation we ourselves report all create strong dependence between columns. Under dependent features, marginal (interventional) Shapley estimators evaluate the model off the data manifold and can distribute credit in ways that do not reflect the conditional structure of the data \cite{aas2021explaining}. Our use of a 100-sample interventional background in \texttt{TreeExplainer} is precisely this setting, and the background is small. More fundamentally, Shapley values are a mathematical allocation, not a causal or human-facing explanation; the properties that make them attractive do not by themselves license the inferences users tend to draw from them, and recovering them for causal purposes requires assumptions the method does not supply \cite{kumar2020problems}. Two consequences follow for this work. Attribution split between two tightly linked positions should not be interpreted as evidence about which is functional; and our restriction of the SHAP analysis to the top 50 features by impurity importance imposes a prior selection step whose effect on the resulting attributions we did not quantify.

The recovery of \emph{rpoB} S450L, \emph{katG} S315T and \emph{embB} M306 should be read in that light. It is a necessary sanity check, and it is reassuring that the pipeline is not learning something spurious -- but it is weak validation. These are the most prevalent resistance mutations globally \cite{who2023catalogue}, so any competent model trained on this data would surface them, and recovering known biology demonstrates the absence of an obvious failure rather than the presence of new insight.

\subsection*{Cross-drug attribution and population structure}

The cross-drug analysis shows that resistance-associated features for one drug (e.g. \emph{katG} S315T) retain measurable attribution in a different drug's model, with no known mechanistic link. The most parsimonious explanation is population-level linkage between resistance mutations in MDR/XDR lineages rather than a causal cross-drug effect. This is a recognized hazard for machine learning on this organism specifically: the \emph{M. tuberculosis} complex has a highly clonal population structure, with negligible horizontal gene transfer and low recombination, so isolates are not statistically independent and models can learn lineage markers that travel with resistance rather than the determinants themselves \cite{billows2023feature}. Our pooled, lineage-agnostic design does nothing to guard against this. We therefore present the cross-drug attribution shifts as a measurement of correlation structure in the training data, not as a biological finding; on the evidence available they are as consistent with confounding as with anything else, and we would discourage readers from reading them as novel resistance biology. Methods for the problem already exist -- stratification, lineage-aware feature selection, and feature-weighted random forests have all been evaluated for exactly this task \cite{billows2023feature} -- and applying one is the single most important next step for this pipeline.

\subsection*{Class imbalance and the value of a susceptible call}

Class imbalance is substantial for two of the four drugs (isoniazid 86.0\% resistant; rifampicin 74.2\% resistant), reflecting a cohort assembled from resistance-focused studies rather than a representative clinical population. While \texttt{class\_weight="balanced"} mitigates its effect during training, headline accuracy and AUC-ROC on an imbalanced test set obscure where the model actually fails: isoniazid susceptible-class precision is 0.61 against 0.98 for the resistant class (Table~\ref{tab:perdrug}). In clinical terms the more consequential error -- calling a resistant isolate susceptible, and thereby endorsing an ineffective regimen -- is the one the imbalance makes hardest to characterize reliably, because susceptible isolates are scarce in the test split. A prospective, consecutively sampled cohort would be needed to estimate that error rate honestly.

\subsection*{Relation to curated catalogues}

Relative to rule-based catalogues such as the WHO catalogue of mutations, now in its second edition and covering more than 30{,}000 variants derived from over 52{,}000 isolates \cite{who2023catalogue}, this pipeline trades curated certainty for (i) native support for four drugs from a single unified pipeline, (ii) a continuous, per-isolate quantitative explanation rather than a binary present/absent mutation flag, and (iii) the ability to surface signal the catalogue does not encode without waiting for a re-curation cycle. The comparison should not be read as a contest. Curated catalogues carry expert adjudication and a documented evidence grade per variant that a learned model has no equivalent of, and independent evaluation on regional cohorts is how their real diagnostic performance has been established \cite{garciamarin2024role}. A catalogue and a model are complementary: the catalogue is the reference standard against which an attribution should be checked, and unexplained attribution is a lead to investigate rather than a finding in itself.

\subsection*{Limitations}

Beyond the issues above, several constraints bound this work. The feature set is restricted to single-nucleotide variants within nine known AMR genes, so indels and structural variants -- which contribute substantially to loss of \emph{pncA} function \cite{yadon2017pnca} -- are excluded entirely, and restricting to known genes forecloses discovery of determinants outside them. Evaluation used a single train/test split per benchmarking run rather than repeated splits or a held-out external cohort, so the reported differences between models carry no confidence intervals and generalization to under-represented regions and lineages is untested. Only four of the 23 drugs present in the underlying phenotype table were modelled. The reference-matching encoding treats any position absent from a sample's VCF as wild type, which conflates true reference calls with positions of low coverage and will systematically underestimate variant burden in poorly sequenced isolates. Finally, the demonstration isolate used for the case study is a single MDR sample and illustrates the output format rather than validating it.

\subsection*{Future directions}

These limitations suggest a clear order of priority. Lineage-aware modelling comes first, since it determines whether the cross-drug findings reported here are biological or artefactual, and established methods for this dataset type are available to adopt rather than invent \cite{billows2023feature}. Second, the interpretability layer should be re-derived with estimators that respect feature dependence \cite{aas2021explaining}, and the effect of the top-50 pre-selection quantified, before any attribution is offered as evidence about a specific locus. Third, evaluation should move to an external cohort with documented phenotyping -- the CRyPTIC compendium \cite{cryptic2022compendium} offers quantitative MICs and would allow the pyrazinamide label-noise question to be separated from the representation question -- and should report against the same isolates used by existing tools so that a genuine comparison with TB-Profiler, GenTB and DeepAMR becomes possible \cite{phelan2019tbprofiler,groschel2021gentb,yang2019deepamr,ngo2019benchmarking}. Fourth, the representation should move beyond individual nucleotide positions toward variant- or protein-level features that can aggregate rare loss-of-function events, which is where the pyrazinamide gains are most likely to come from. Extension to second-line drugs and richer feature types follows from these.

Taken together, this work demonstrates that an interpretable, multi-drug pipeline can be assembled end to end from public sequencing data and open-source tooling, and released in full. Its value lies in that reproducibility and in the per-isolate explanation layer, not in predictive performance, which trails established tools. We offer FORUM-TB as a transparent baseline and a shared resource against which more sophisticated, lineage-aware and clinically validated interpretable models can be developed and -- given a common benchmark -- fairly compared.
